\documentclass[11pt,a4paper]{article}
\usepackage[a4paper,margin=2.1cm]{geometry}
\usepackage[spanish,es-noquoting]{babel}
\usepackage[T1]{fontenc}
\usepackage[utf8]{inputenc}
\usepackage{lmodern,microtype}
\usepackage{amsmath,amssymb,amsthm,mathtools}
\usepackage{booktabs,array,longtable,enumitem}
\usepackage{xcolor,hyperref}
\hypersetup{colorlinks=true,linkcolor=blue!55!black,urlcolor=blue!55!black,citecolor=blue!55!black}
\setlength{\parindent}{0pt}
\setlength{\parskip}{0.65em}
\definecolor{HMTblue}{RGB}{18,70,140}
\definecolor{HMTred}{RGB}{150,35,35}
\definecolor{HMTgreen}{RGB}{20,115,75}
\definecolor{HMTgold}{RGB}{150,105,15}
\definecolor{HMTsoft}{RGB}{245,248,255}
\newcommand{\blue}[1]{\textcolor{HMTblue}{#1}}
\newcommand{\red}[1]{\textcolor{HMTred}{#1}}
\newcommand{\green}[1]{\textcolor{HMTgreen}{#1}}
\newcommand{\gold}[1]{\textcolor{HMTgold}{#1}}
\newcommand{\F}{\mathbb F}
\newcommand{\Z}{\mathbb Z}
\newcommand{\Pone}{\mathbb P^1}
\newcommand{\APP}{\operatorname{APP}}
\newcommand{\HMT}{\mathrm{HMT}}
\newcommand{\AW}{A_W}
\newcommand{\CW}{C_W}
\newcommand{\Wdoce}{W_{12}}
\newcommand{\Vonce}{V_{11}}
\newcommand{\one}{\mathbf 1}
\newtheorem{definicion}{Definición}
\newtheorem{lema}{Lema}
\newtheorem{proposicion}{Proposición}
\newtheorem{teorema}{Teorema}
\newtheorem{observacion}{Observación}
\newtheorem{conjetura}{Conjetura}
\title{\bfseries \textcolor{HMTblue}{Puerta Hensel--Paley--Witt}\
\large Derivación de la polaridad \(A_W\), cierre Golay--Witt y puente de once modos}
\author{Ventana técnica HMT}
\date{}
\begin{document}
\maketitle
\begin{abstract}
Se cierra la puerta conceptual pendiente del macroobjeto HMT. La dualidad \(A_W\) no debe presentarse como una matriz insertada a mano, sino como la polaridad Paley--Witt inducida por la frontera reversible de la célula Hensel. La cadena es
\[
\blue{\APP+U_6+\text{Hensel}\longrightarrow U_9\longrightarrow \Pone(\F_5)\longrightarrow \chi_5\longrightarrow A_W\longrightarrow W_{12}.}
\]
Se demuestra la ortogonalidad por sumas de caracteres cuadráticos en \(\F_5\), se verifica el enumerador de Golay ternario y se explica cómo esta misma estructura prepara el módulo de once modos
\(\mathbb R^{12}=\mathbb R\one\oplus V_{11}\), con \(V_{11}\cong3\oplus3'\oplus5\) en la lectura icosaédrica.
\end{abstract}
\tableofcontents

\section{Dictamen}
La formulación correcta ya no es:
\[
\text{``falta derivar }A_W\text{ desde APP/cubo/carry''.}
\]
La formulación correcta es:
\[
\boxed{\blue{\text{hay que auditar la cadena }U_9\to \Pone(\F_5)\to\chi_5\to A_W.}}
\]
Hecha esa auditoría, \(A_W\) deja de ser un ansatz libre. Se obtiene como polaridad Paley de una línea proyectiva de seis puntos.

\section{De APP/Hensel a la frontera reversible \(U_9\)}
El núcleo Hensel es
\[
\beta:\F_3^2\to\Z/9\Z,
\qquad
\beta(a,b)=a+3b\pmod9.
\]
Así:
\[
U_6=\F_3^6\xrightarrow{\beta^3}(\Z/9\Z)^3.
\]
APP es una carta de dos coordenadas de esta célula cúbico--Hensel. La parte reversible de una coordenada \(\Z/9\Z\) no son los nueve estados, sino sus unidades:
\[
\boxed{\green{U_9=(\Z/9\Z)^\times=\{1,2,4,5,7,8\}.}}
\]
Los estados \(0,3,6\) son singulares respecto del producto. Por tanto, si buscamos una frontera orientada e invertible de la célula APP/Hensel, la candidata natural es \(U_9\), no todo \(\Z/9\Z\).

\begin{observacion}
Esta es la primera reducción no arbitraria: APP separa unidades y no-unidades. La frontera reversible de Hensel tiene seis puntos.
\end{observacion}

\section{De \(U_9\) a \(\Pone(\F_5)\)}
El grupo de unidades \(U_9\) es cíclico de orden seis. Un generador es \(2\):
\[
1,\;2,\;4,\;8,\;7,\;5,\;1.
\]
Una geometría proyectiva natural con seis puntos es
\[
\Pone(\F_5)=\{\infty,0,1,2,3,4\}.
\]
La identificación HMT preserva el ciclo de unidades mediante el orden
\[
\boxed{\gold{(\infty,0,1,4,2,3).}}
\]
Concretamente:
\[
1\mapsto\infty,
\quad
2\mapsto0,
\quad
4\mapsto1,
\quad
8\mapsto4,
\quad
7\mapsto2,
\quad
5\mapsto3.
\]
Este orden es el ciclo de la transformación de Möbius
\[
T(x)=\frac{1}{3x+1}\quad \text{en }\F_5,
\]
con
\[
\infty\mapsto0\mapsto1\mapsto4\mapsto2\mapsto3\mapsto\infty.
\]

\begin{proposicion}[Projectivización de la frontera reversible]
La frontera reversible \(U_9\) de la célula Hensel admite una projectivización natural como \(\Pone(\F_5)\) al preservar su ciclo de unidades.
\end{proposicion}

\begin{observacion}
Aquí nace el componente pentádico: seis unidades Hensel se leen como \(1+5\), un punto infinito más cinco puntos finitos. El 5 no entra por gusto; entra como parte finita de la frontera proyectiva.
\end{observacion}

\section{El carácter cuadrático \(\chi_5\)}
El carácter cuadrático de \(\F_5\) es
\[
\chi_5(0)=0,
\qquad
\chi_5(1)=\chi_5(4)=+1,
\qquad
\chi_5(2)=\chi_5(3)=-1.
\]
Esta es la polaridad mínima disponible en la línea proyectiva: distingue residuos y no-residuos cuadráticos.

\section{Construcción de la matriz Paley--Witt}
Indexamos por
\[
(\infty,0,1,4,2,3).
\]
Definimos la matriz signada \(C\) por:
\[
C_{xx}=0,
\]
\[
C_{\infty,y}=C_{y,\infty}=1,
\]
\[
C_{x,y}=\chi_5(x-y)\quad(x,y\in\F_5,
\ x\ne y).
\]
En ese orden:
\[
C=
\begin{pmatrix}
0&1&1&1&1&1\\
1&0&1&1&-1&-1\\
1&1&0&-1&1&-1\\
1&1&-1&0&-1&1\\
1&-1&1&-1&0&1\\
1&-1&-1&1&1&0
\end{pmatrix}.
\]
Reduciendo \(-1\equiv2\pmod3\), obtenemos:
\[
A_W=
\begin{pmatrix}
0&1&1&1&1&1\\
1&0&1&1&2&2\\
1&1&0&2&1&2\\
1&1&2&0&2&1\\
1&2&1&2&0&1\\
1&2&2&1&1&0
\end{pmatrix}\pmod3.
\]
Esta forma es gauge/monomialmente equivalente a la normalización usada en la auditoría con \(A_WA_W^T\equiv2I\pmod3\). Ambas representan la misma polaridad Witt.

\section{Ortogonalidad por suma de caracteres}
\begin{teorema}[Ortogonalidad Paley]
La matriz signada \(C\) satisface
\[
\boxed{C C^T=5I.}
\]
\end{teorema}

\begin{proof}
La norma de cada fila es 5, porque cada fila contiene cinco entradas no nulas de valor \(\pm1\).

Sea \(x\in\F_5\). El producto escalar de la fila \(\infty\) con la fila \(x\) es
\[
\sum_{z\in\F_5}\chi_5(x-z)=0,
\]
porque \(\chi_5\) es un carácter no trivial.

Para \(x\ne y\) en \(\F_5\), el producto escalar entre las filas finitas es
\[
1+
\sum_{z\in\F_5}\chi_5(x-z)\chi_5(y-z).
\]
Usamos
\[
\sum_{z\in\F_5}\chi_5((x-z)(y-z))=-1
\quad(x\ne y),
\]
propiedad estándar de las sumas de caracteres cuadráticos. Por tanto el producto escalar es
\[
1-1=0.
\]
Así todas las filas son ortogonales y tienen norma 5.
\end{proof}

Reduciendo módulo 3:
\[
5\equiv2\pmod3,
\]
por tanto:
\[
\boxed{A_WA_W^T\equiv2I\pmod3.}
\]

\section{Código gráfico y autodualidad}
Definimos
\[
C_W=\{(q,qA_W):q\in\F_3^6\}.
\]
Con
\[
G=[I_6\mid A_W],
\]
se tiene
\[
GG^T=I_6+A_WA_W^T\equiv I_6+2I_6\equiv0\pmod3.
\]
Así el código es auto-ortogonal. Como tiene dimensión 6 en longitud 12, es autodual.

La enumeración produce:
\[
\boxed{1+264y^6+440y^9+24y^{12}.}
\]
Los 264 codewords de peso 6 dan 132 soportes bajo \(c\sim-c\). Esos soportes son las hexadas de \(W_{12}\).

\section{Witt pequeño}
Los números de incidencia son:
\[
\lambda_0=132,
\quad
\lambda_1=66,
\quad
\lambda_2=30,
\quad
\lambda_3=12,
\quad
\lambda_4=4,
\quad
\lambda_5=1.
\]
En particular:
\[
\boxed{\lambda_5=1.}
\]
Esto significa:
\[
\boxed{\text{todo 5-subconjunto está contenido en una única hexada}.}
\]
Por tanto:
\[
\boxed{W_{12}=S(5,6,12).}
\]

\section{Cierre de la puerta conceptual}
La puerta conceptual pendiente eran tres decisiones:
\begin{enumerate}[label=\textbf{\arabic*.}]
\item Tomar las unidades \(U_9\) como frontera reversible.
\item Projectivizarlas como \(\Pone(\F_5)\) preservando el ciclo de unidades.
\item Tomar \(\chi_5\) como polaridad cuadrática de esa frontera.
\end{enumerate}
Bajo esas tres decisiones, que son naturales en la gramática Hensel/APP, se obtiene \(A_W\). Luego Golay/Witt cae por ortogonalidad Paley.

\[
\boxed{\APP+\text{Hensel}\to U_9\to\Pone(\F_5)\to\chi_5\to A_W\to W_{12}.}
\]

\section{Puente de once modos}
La misma rueda dodecafásica permite otra lectura. Sea
\[
V_{12}=\mathbb R^{12},
\qquad
\one=(1,\ldots,1).
\]
El modo global es \(\mathbb R\one\). El espacio activo es
\[
V_{11}=\one^\perp=
\{x\in\mathbb R^{12}:\sum_i x_i=0\}.
\]
Entonces
\[
\boxed{V_{12}=\mathbb R\one\oplus V_{11}.}
\]
Si los 12 puntos se leen con la acción icosaédrica \(A_5\), la representación de permutación se descompone como
\[
\boxed{\mathbb R^{12}\cong1\oplus3\oplus3'\oplus5.}
\]
Por tanto:
\[
\boxed{V_{11}\cong3\oplus3'\oplus5=3\oplus8.}
\]
Lectura HMT:
\[
3=\text{sector visible/cúbico},
\]
\[
3'=\text{sector dual/contracara},
\]
\[
5=\text{sector pentádico de cierre},
\]
\[
8=3'\oplus5=\text{sector interno octádico}.
\]
Esto no debe venderse como igualdad literal con M-teoría. La frase correcta es:
\[
\boxed{\text{HMT no reproduce la teoría M; explica por qué una física de frontera puede ver once modos activos.}}
\]

\section{Conclusión}
La cadena fuerte queda:
\[
\boxed{\text{célula Hensel}\to\text{frontera reversible}\to\text{línea proyectiva}\to\text{Paley}\to\text{Witt}.}
\]
Y el macroobjeto HMT queda formulado como una célula cúbico--Hensel--Witt:
\[
\boxed{
\mathcal X_{\rm HMT}=(U_6,\beta^3,\partial_{\APP},\kappa,A_W,C_W,W_{12},\mathcal G_9,\mathcal R_{1000}).
}
\]
La frase final:
\[
\boxed{\text{la polaridad Witt emerge de la frontera Hensel de APP; no se inserta como matriz arbitraria.}}
\]
\end{document}
